\documentclass[11pt]{article}
% Uses the canonical latex-docs semantic style in-repository. The local
% lightweight fallback is for independent compilation of this distribution.
\IfFileExists{technical-user-manual.sty}{\usepackage{technical-user-manual}}{\usepackage{../support/aaoa-standalone}}
\usepackage{mathtools}
\setDocTitle{Chapter 7: Permutations}
\setDocSubtitle{Guided reading, derivations and original exercises}
\setDocVersion{1.0}
\setDocStatus{Educational reference / unapproved working material}
\setDocOwner{Algorithm Engineering}
\setDocAuthor{Original companion material}
\setDocSummary{Coverage of primary book Chapter 7 with independently authored examples, proof obligations and lab transfer.}
\providecommand{\Expect}{\mathbb{E}}
\providecommand{\Prob}{\mathbb{P}}
\providecommand{\Var}{\operatorname{Var}}
\providecommand{\bigoh}{\mathcal{O}}
\begin{document}
\makedoctocpage
\section{Reading map and prerequisite check}
\noindent\textbf{Primary source sections:} 7.1 Basic Properties; 7.2 Algorithms; 7.3 Representations; 7.4 Enumeration; 7.5 CGFs; 7.6 Inversions and Insertion Sorts; 7.7 Left-to-Right Minima and Selection Sort; 7.8 Cycles and In Situ Permutation; 7.9 Extremal Parameters (printed pp. 347--406).\par
\noindent\textbf{Prerequisites:} Completed preceding mathematical-method chapters or equivalent demonstrated mastery.\par
\section{Learning outcomes}
\begin{enumerate}
\item Derive expected inversion count with indicator variables.
\item Connect a permutation statistic to an exact implementation count.
\item Recognize cycles and records as alternative statistics with different costs.
\end{enumerate}
\section{Guided technical analysis}
\subsection{Pair indicators}
For a uniformly random permutation, let $I_{ij}=1$ if $i<j$ yet $a_i>a_j$. Symmetry gives $\Expect[I_{ij}]=1/2$, hence
\[
\Expect\!\left[\sum_{i<j}I_{ij}\right]=\frac{n(n-1)}{4}.
\]
No independence assumption is needed to sum expectations.
\subsection{Insertion sort operation identity}
For insertion sort implemented by repeatedly shifting a larger predecessor right by one, the number of shifts equals the input inversion count. This is exact for that specific implementation, not necessarily for all swap-based or binary-insertion variants.
\subsection{Permutation cycles}
Write a permutation as disjoint cycles. Cycle count differs from inversion count and depends on the mapping representation; an in-place permutation routine may use different moves per cycle. Identify fixed points, cycles and representation before deriving its cost.
\section{Independent problem set}
\begin{enumerate}
\item Enumerate all $4!$ permutations; calculate inversion distribution and its exact mean.
\item Prove insertion-sort shifts equal inversions using an invariant on sorted prefixes.
\item Compare the same permutation encoded as an array and as disjoint cycles; count visits for in-place rearrangement.
\end{enumerate}
\section{Laboratory transfer}
\textbf{Required laboratory:} LAB-CPP02: enumerate and instrument insertion-sort shift counts.\par
\noindent\textbf{Minimum recorded evidence:} repository commit or source hash; test inputs/seeds; actual operation counts; predicted values; one explicit model limitation.\par
\section{Mastery gate and self-review}
Test all permutations of sizes 1--7 (or justify a smaller exhaustive bound) with a correct inversion oracle.\par
\begin{enumerate}\item Can I define the random object/input exactly? \item Is my identity exact, asymptotic, or empirical? \item Which base cases and counterexamples did I check? \item What would falsify my performance hypothesis?\end{enumerate}
\section{Suggested notebook structure}
Use five named sections: model and invariants; derivation; deterministic small instances; randomized replication; discrepancy and revised hypothesis. Keep a proof separate from measurement results.

\end{document}
